% ME16020C
% Hypoglykämie
% 14.0
% 2013-11-30

\label{m:hypoglykaemie}Vitale
Bedrohung je nach Bewusstseinsgrad

\EE{Taktik}: \Speech{Zucker zuführen}
\begin{itemize}
  \item Vitale Bedrohung einschätzen. \TxBeiVit \TxMassVitMK
  
  \item BZ messen! Zur Diagnosestellung bzw. \E{vor} und \E{nach}
  jeder Zuckergabe
  \item Zucker zuführen: Wenn der Patient bewusstseinsklar
  ist, Traubenzucker oder Zuckerwasser geben. Danach \enquote{feste
  Speisen} nachessen lassen, da der Traubenzucker nur eine
  kurze Wirkdauer hat.
  
  Wenn der Patient nicht mehr bewusstseinsklar ist, kann die
  Zuckergabe nur über eine venöse Gabe, z.\,B. durch einen
  Notarzt oder NKV, erfolgen (getrübte Patienten dürfen
  aufgrund der Aspirationsgefahr nichts mehr essen).
\end{itemize}

\Customized{
\begin{description}
  \item[\CustAsboe] Bei der Hypoglykämie mit Bewusstseinsstörung ist die Gabe
  von Glukose i.\,v. und kristalloider Infusionslösung gem. Algorithmus durch NKV vorgesehen.
  \cite{Asb921Memo-AL-2011-000001}.
\end{description}
}